\documentclass[12pt]{article}
\usepackage[T1]{fontenc}
\usepackage[utf8]{inputenc}
\usepackage{lmodern}
\usepackage{textcomp}
\usepackage{enumitem}
\usepackage{geometry}
\geometry{margin=1in}
\begin{document}
\begin{center}
Answer Key - Economics - Graduate Level Assessment 21 \
\small (Answers are ordered exactly as in the question paper.)
\end{center}

\section*{Multiple Choice}
\begin{enumerate}[label=\arabic*.]
\item \textbf{Answer:} A
\\ \textit{Options:} A) An increase in unemployment rates; B) A decrease in government spending; C) A decrease in the supply of money; D) An increase in savings; E) An increase in interest rates
\\ \textit{Note:} An increase in unemployment rates can lead to increased government intervention through fiscal policies, such as stimulus spending, which can boost consumer confidence and aggregate demand, thus shifting the aggregate demand curve to the right.
\item \textbf{Answer:} D
\\ \textit{Options:} A) No, it is a principle that applies only to consumer choice in the market.; B) No, it is only applicable to the exchange of services, not goods.; C) Yes, it specifically applies to trade agreements between neighboring countries.; D) No, it applies wherever productive units have different relative efficiencies.; E) Yes, it only applies to international trade.
\\ \textit{Note:} The principle of comparative advantage applies not only to international trade but also to any scenario where different productive units can achieve varying levels of efficiency in producing goods or services, leading to optimal resource allocation and specialization regardless of the context.
\item \textbf{Answer:} C
\\ \textit{Options:} A) Linearly increasing; B) Inverted U-shape; C) Horizontal; D) Upward sloping; E) Vertical
\\ \textit{Note:} A horizontal demand curve indicates that Matt's total utility from consuming bratwurst increases at a constant rate regardless of the quantity consumed, suggesting he is willing to purchase any amount at a constant price without diminishing marginal utility.
\item \textbf{Answer:} A
\\ \textit{Options:} A) not found by adding product demand curves horizontally or vertically; B) a vertical summation of market demand curves; C) found by adding product demand curves both horizontally and vertically; D) a horizontal subtraction of firm demand curves; E) a simple aggregation of demand curves for individual goods
\\ \textit{Note:} The aggregate demand curve represents the total demand for goods and services in an economy at various price levels and is derived from the overall relationships between price levels and total output, rather than simply summing individual product demand curves either horizontally or vertically.
\item \textbf{Answer:} A
\\ \textit{Options:} A) Chicago school, Stiglitz; B) Marxist school, Minsky; C) Behavioral school, Akerlof; D) Neoclassical school, Friedman; E) Neoclassical school, Coase
\\ \textit{Note:} The Chicago school, known for its emphasis on free markets and supply-side economics, defends Say's Law, which posits that supply creates its own demand, while Joseph Stiglitz, a prominent economist, disputes this theory by highlighting the role of information asymmetries and market failures in economic fluctuations.
\end{enumerate}

\section*{True / False}
\begin{enumerate}[label=\arabic*.]
\setcounter{enumi}{5}
\item \textbf{Answer:} True
\\ \textit{Options:} A) True; B) False
\\ \textit{Note:} A recovery from a recession typically involves a resurgence in economic activity, reflected in indicators such as GDP growth, employment rates, and consumer confidence, which collectively signal a return to stability and positive economic momentum.
\item \textbf{Answer:} False
\\ \textit{Options:} A) True; B) False
\\ \textit{Note:} The statement is false because the actual growth in real GDP per capita during that period was significantly lower than 50 percent, reflecting a more moderate economic expansion than suggested.
\item \textbf{Answer:} False
\\ \textit{Options:} A) True; B) False
\\ \textit{Note:} The statement is false because a perfectly competitive firm maximizes profit by continuing to operate as long as the price exceeds the marginal cost, not merely the maximum average total cost.
\item \textbf{Answer:} False
\\ \textit{Options:} A) True; B) False
\\ \textit{Note:} Classical analysis holds that during a recession, contractionary monetary policy, which reduces the money supply and increases interest rates, would further decrease aggregate demand, worsening economic conditions rather than facilitating recovery.
\item \textbf{Answer:} False
\\ \textit{Options:} A) True; B) False
\\ \textit{Note:} The statement is false because marginal cost intersects average variable cost at the minimum point of average variable cost, where average variable cost is equal to marginal cost, not greater.
\end{enumerate}

\section*{Long Form Response}
\begin{enumerate}[label=\arabic*.]
\setcounter{enumi}{10}
\item \textbf{Answer:} Money serves as a unit of account by providing a consistent measure for pricing goods and services, which facilitates economic transactions and comparisons. For instance, the price of a 20-ounce bottle of pop at \$1.25 allows consumers to gauge its value relative to other products, streamlining decision-making processes in purchasing. This uniform pricing structure enables individuals and businesses to assess costs, budget allocations, and the overall value of transactions, fostering efficient market operations. Moreover, it simplifies the comparison of prices across different markets and timeframes, ensuring that consumers can make informed choices based on relative value rather than subjective assessments. Ultimately, as a unit of account, money enhances clarity and efficiency in economic interactions.
\\ \textit{Note:} correct because it illustrates how money as a unit of account standardizes the pricing of goods, like a 20-ounce bottle of pop at \$1.25, enabling consumers to make informed comparisons and decisions, thereby enhancing the efficiency of economic transactions.
\item \textbf{Answer:} Marginal cost (MC) equals average variable cost (AVC) and average total cost (ATC) at their minimum points, where the curves intersect. This occurs because, when MC is below AVC or ATC, the average costs are decreasing, and when MC is above, they are increasing; thus, MC intersects AVC and ATC at their average points. The implications of this relationship for production decisions are significant; firms should aim to produce at levels where MC equals AVC and ATC to minimize costs and maximize efficiency. If a firm's production level is such that MC exceeds AVC or ATC, it indicates inefficiency and may warrant a reevaluation of production strategies. Conversely, operating at levels where MC is below these averages suggests potential for increased output without incurring higher average costs.
\\ \textit{Note:} correct because it accurately describes that marginal cost intersects average variable cost and average total cost at their minimum points, reflecting the relationship between costs and production efficiency, which is crucial for firms to minimize costs and optimize production decisions.
\end{enumerate}

\vfill
\noindent\textit{End of Answer Key}
\end{document}